\documentclass[12pt,a4paper]{article}
\usepackage[margin=2.5cm]{geometry}
\usepackage{listings}
\usepackage{xcolor}
\usepackage{amsmath}
\usepackage{booktabs}
\usepackage{hyperref}
\usepackage{graphicx}
\usepackage{float}

% XML/code listing style
\definecolor{codebg}{rgb}{0.95,0.95,0.95}
\definecolor{codegreen}{rgb}{0,0.5,0}
\definecolor{codepurple}{rgb}{0.5,0,0.5}
\definecolor{codeblue}{rgb}{0,0,0.7}

\lstdefinestyle{xml}{
    backgroundcolor=\color{codebg},
    basicstyle=\ttfamily\small,
    keywordstyle=\color{codeblue},
    stringstyle=\color{codegreen},
    commentstyle=\color{gray}\itshape,
    breaklines=true,
    frame=single,
    numbers=left,
    numberstyle=\tiny\color{gray},
    tabsize=4
}

\lstdefinestyle{python}{
    backgroundcolor=\color{codebg},
    basicstyle=\ttfamily\small,
    keywordstyle=\color{codeblue}\bfseries,
    stringstyle=\color{codegreen},
    commentstyle=\color{gray}\itshape,
    breaklines=true,
    frame=single,
    numbers=left,
    numberstyle=\tiny\color{gray},
    tabsize=4,
    language=Python
}

\title{\textbf{Furuta Pendulum: SolidWorks to MuJoCo}\\
\large Build Documentation}
\author{}
\date{\today}

\begin{document}

\maketitle
\tableofcontents
\newpage

%─────────────────────────────────────────────────────────────────────────────
\section{System Overview}

The Furuta (rotary inverted) pendulum consists of three rigid bodies:

\begin{itemize}
    \item \textbf{Fixed motor} -- motor housing, bracket, and bearings, fixed to the world
    \item \textbf{Rotary arm} -- horizontal arm driven by the motor shaft (1 actuated DOF)
    \item \textbf{Pendulum rod} -- rod with tip mass, free to swing from the end of the arm (1 passive DOF)
\end{itemize}

\noindent The system has two joints:
\begin{itemize}
    \item \textbf{Shoulder} ($\phi$): arm rotation about the motor shaft axis
    \item \textbf{Elbow} ($\theta$): pendulum rotation about the pivot at the arm tip
\end{itemize}

\noindent Motor: \textbf{25GA20E260}, 7.4V DC, 1:20 gear ratio.\\
Pendulum angle sensor: \textbf{AS5600} magnetic encoder (12-bit).

%─────────────────────────────────────────────────────────────────────────────
\section{STL Export from SolidWorks}

Each moving body is exported as a separate STL file from the SolidWorks assembly.
The key requirement is that \textbf{all meshes must be exported in the same global coordinate frame} -- this ensures the relative positions between parts are already encoded in the mesh vertex coordinates.

\subsection*{Files}
\begin{center}
\begin{tabular}{lll}
\toprule
\textbf{File} & \textbf{Body} & \textbf{Description} \\
\midrule
\texttt{fixed\_motor.stl} & fixed\_motor & Motor housing, bracket, bearings \\
\texttt{arm\_2.STL}       & arm          & Updated rotary arm + flange coupler \\
\texttt{rod\_2.STL}       & pendulum\_rod & Updated pendulum rod with tip weight \\
\bottomrule
\end{tabular}
\end{center}

\subsection*{Coordinate System}
SolidWorks uses \textbf{X-up} as the vertical axis, while MuJoCo uses \textbf{Z-up}.
The remapping is applied globally via the \texttt{<default>} geom:

\begin{lstlisting}[style=xml]
<default>
    <geom type="mesh" xyaxes="0 0 1 1 0 0"/>
</default>
\end{lstlisting}

\noindent This defines the mesh local axes in world frame:
\begin{itemize}
    \item Mesh local X $\rightarrow$ World Z (up)
    \item Mesh local Y $\rightarrow$ World X
    \item Mesh local Z $\rightarrow$ World Y (derived as cross product)
\end{itemize}

%─────────────────────────────────────────────────────────────────────────────
\section{Mesh Re-centering with \texttt{refpos}}

SolidWorks exports meshes at their global assembly coordinates.
The motor center in the SolidWorks assembly is at $(602.303,\ 877.409,\ 1307.925)$ mm.
To re-center all meshes at the body origin, \texttt{refpos} subtracts this offset from all vertex positions:

\begin{lstlisting}[style=xml]
<asset>
    <mesh name="fixed_motor" file="cad/fixed_motor.stl"
          scale="0.001 0.001 0.001"
          refpos="602.303 877.409 1307.925"/>
    <mesh name="arm" file="cad/arm_2.STL"
          scale="0.001 0.001 0.001"
          refpos="602.303 877.409 1307.925"/>
    <mesh name="pendulum_rod" file="cad/rod_2.STL"
          scale="0.001 0.001 0.001"
          refpos="602.303 877.409 1307.925"/>
</asset>
\end{lstlisting}

\noindent Notes:
\begin{itemize}
    \item \texttt{scale="0.001 0.001 0.001"} converts SolidWorks mm $\rightarrow$ MuJoCo meters
    \item All three meshes use the same \texttt{refpos} so they remain correctly assembled
\end{itemize}

%─────────────────────────────────────────────────────────────────────────────
\section{Kinematic Structure}

The bodies are nested to reflect the kinematic chain:

\begin{lstlisting}[style=xml]
<body name="fixed_motor" pos="0 0 0.15">        <!-- fixed to world -->
    <body name="arm">
        <joint name="shoulder" .../>             <!-- DOF 1 -->
        <body name="pendulum_rod">
            <joint name="elbow" .../>            <!-- DOF 2 -->
        </body>
    </body>
</body>
\end{lstlisting}

\noindent The \texttt{fixed\_motor} body is placed at $z = 0.15$ m above the world origin.
Since it has no joint, it is rigidly fixed to the world.

%─────────────────────────────────────────────────────────────────────────────
\section{Joint Definitions}

\subsection{Shoulder Joint}

The shoulder joint allows the arm to rotate about the motor shaft axis.
The motor shaft center was measured in SolidWorks at $(628.05,\ 877.41,\ 1307.92)$ mm.
After \texttt{refpos} subtraction and coordinate remapping, the joint anchor is:

\begin{equation}
\text{pos} = (0,\ 0,\ 0.02575) \text{ m}
\end{equation}

\noindent The shaft is aligned with the world Z-axis (vertical axis in SolidWorks $\rightarrow$ world Z in MuJoCo):

\begin{lstlisting}[style=xml]
<joint name="shoulder" pos="0 0 0.02575" axis="0 0 1" damping="0.01"
       limited="true" range="-135 135" solreflimit="0.005 1"/>
\end{lstlisting}

\subsection{Elbow Joint}

The elbow joint allows the pendulum to swing freely from the arm tip.  After
switching from the original \texttt{arm}/\texttt{rod} CAD to
\texttt{arm\_2}/\texttt{rod\_2}, the pivot and axis were updated to match the
new assembly.  The current convention uses \texttt{ref="0"}:

\begin{lstlisting}[style=xml]
<joint name="elbow" pos="0 -0.049505 0.056747" axis="0 1 0" ref="0"/>
\end{lstlisting}

\noindent With this convention:
\begin{itemize}
    \item $\theta = 0$ rad: pendulum upright (unstable equilibrium / balance target)
    \item $\theta = \pm\pi$ rad: pendulum hanging down (stable equilibrium)
\end{itemize}

%─────────────────────────────────────────────────────────────────────────────
\section{Physical Properties}

\subsection{Mass}

Masses were measured on a scale and set directly on each \texttt{<geom>}:

\begin{center}
\begin{tabular}{lrr}
\toprule
\textbf{Body} & \textbf{Measured mass (g)} & \textbf{XML value (kg)} \\
\midrule
fixed\_motor  & 45.60 & 0.04560 \\
arm           & 25.29 & 0.02529 \\
pendulum\_rod & 18.88 & 0.01888 \\
\bottomrule
\end{tabular}
\end{center}

\noindent The listed values are the current XML masses for the updated
\texttt{arm\_2.STL} and \texttt{rod\_2.STL} geometry.

\subsection{Inertia}

Inertia tensors are computed automatically by MuJoCo from the mesh geometry and the assigned mass.
No manual inertia specification is required.

%─────────────────────────────────────────────────────────────────────────────
\section{Actuator}

The motor is modeled as a direct torque actuator on the shoulder joint.
Torque limits are derived from the motor datasheet (25GA20E260):

\begin{center}
\begin{tabular}{lrl}
\toprule
\textbf{Parameter} & \textbf{Value} & \textbf{Source} \\
\midrule
Rated torque  & 0.039 Nm & $0.4 \text{ kg.cm} \times 0.0981$ \\
Stall torque  & 0.147 Nm & $1.5 \text{ kg.cm} \times 0.0981$ \\
Gear ratio    & 1:20     & datasheet \\
Gearbox efficiency & 90\% & assumed \\
Effective max torque & 0.132 Nm & $0.147 \times 0.90$ \\
\bottomrule
\end{tabular}
\end{center}

\begin{lstlisting}[style=xml]
<actuator>
    <motor name="shoulder_motor" joint="shoulder"
           gear="0.132" ctrlrange="-1 1"/>
</actuator>
\end{lstlisting}

\noindent Control input is normalized: $u \in [-1, 1]$ maps linearly to
$[-0.132, 0.132]$ Nm through the motor \texttt{gear} value.
Joint damping of 0.01 Nm$\cdot$s/rad approximates gearbox friction.

%─────────────────────────────────────────────────────────────────────────────
\section{Contact Exclusions}

Adjacent bodies in the kinematic chain overlap at the joint locations by design (shaft inside bearing, etc.).
To prevent spurious internal contact forces, neighboring body pairs are excluded:

\begin{lstlisting}[style=xml]
<contact>
    <exclude body1="fixed_motor" body2="arm"/>
    <exclude body1="arm"         body2="pendulum_rod"/>
    <exclude body1="fixed_motor" body2="pendulum_rod"/>
</contact>
\end{lstlisting}

\noindent All three body pairs are currently excluded to prevent CAD overlap
from creating spurious contact forces during training.

%─────────────────────────────────────────────────────────────────────────────
\section{Complete XML Model}

\begin{lstlisting}[style=xml]
<mujoco model="furuta_pendulum">
    <default>
        <!-- remap SolidWorks X-up to MuJoCo Z-up -->
        <geom type="mesh" xyaxes="0 0 1 1 0 0"/>
    </default>

    <contact>
        <exclude body1="fixed_motor" body2="arm"/>
        <exclude body1="arm" body2="pendulum_rod"/>
        <exclude body1="fixed_motor" body2="pendulum_rod"/>
    </contact>

    <visual>
        <headlight diffuse="0.6 0.6 0.6" ambient="0.3 0.3 0.3"
                   specular="0.1 0.1 0.1"/>
        <rgba haze="0.15 0.25 0.35 1"/>
    </visual>

    <asset>
        <mesh name="fixed_motor" file="cad/fixed_motor.stl"
              scale="0.001 0.001 0.001" refpos="602.303 877.409 1307.925"/>
        <mesh name="arm" file="cad/arm_2.STL"
              scale="0.001 0.001 0.001" refpos="602.303 877.409 1307.925"/>
        <mesh name="pendulum_rod" file="cad/rod_2.STL"
              scale="0.001 0.001 0.001" refpos="602.303 877.409 1307.925"/>
    </asset>

    <worldbody>
        <light pos="0 0 1.5" dir="0 0 -1" diffuse="0.8 0.8 0.8"/>
        <light pos="0.5 0.5 1" dir="-1 -1 -1" diffuse="0.4 0.4 0.4"/>

        <body name="fixed_motor" pos="0 0 0.15">
            <geom mesh="fixed_motor" mass="0.04560"/>
            <body name="arm">
                <joint name="shoulder" pos="0 0 0.02575"
                       axis="0 0 1" damping="0.01"
                       limited="true" range="-135 135" solreflimit="0.005 1"/>
                <geom mesh="arm" mass="0.02529"/>
                <body name="pendulum_rod">
                    <joint name="elbow" pos="0 -0.049505 0.056747"
                           axis="0 1 0" ref="0"/>
                    <geom mesh="pendulum_rod" mass="0.01888"/>
                </body>
            </body>
        </body>
    </worldbody>

    <actuator>
        <motor name="shoulder_motor" joint="shoulder"
               gear="0.132" ctrlrange="-1 1"/>
    </actuator>
</mujoco>
\end{lstlisting}

%─────────────────────────────────────────────────────────────────────────────
\section{Sanity Checks}

Three physics sanity checks were performed with the shoulder joint locked
(arm fixed, pendulum swinging freely):

\subsection{Pendulum Period}

For a simple pendulum: $T = 2\pi\sqrt{L/g}$, where $L$ is the distance from pivot to center of mass.

\begin{center}
\begin{tabular}{lrr}
\toprule
\textbf{Parameter} & \textbf{Value} \\
\midrule
Current pivot-to-CoM distance & 3.60 cm \\
Theoretical period $T$ & 0.381 s \\
Simulated period & Re-check after final CAD freeze \\
\bottomrule
\end{tabular}
\end{center}

\noindent The pivot-to-CoM distance is computed from the current MuJoCo body
frame values:
\texttt{rod\_COM=[0,\ -0.035077,\ 0.089760]} m and
\texttt{elbow\_pos=[0,\ -0.049505,\ 0.056747]} m.

\subsection{Visual Inspection}

The pendulum oscillates smoothly around its stable equilibrium with no jitter or instability.
Contact exclusions prevent internal collision forces between adjacent bodies.

\subsection{Energy Conservation}

The current elbow joint has no explicit damping in XML.  Releasing the
pendulum from 40$^\circ$ past equilibrium should therefore be used as a
final sanity check after the CAD and contact exclusions are frozen:
\begin{itemize}
    \item Expected behaviour: smooth swing with no jitter or artificial contact loss
    \item Any measured real-hardware decay should be used to identify pivot friction
\end{itemize}

%─────────────────────────────────────────────────────────────────────────────
\section{Angle Convention}

With the current \texttt{ref="0"} elbow joint, the angle convention is:

\begin{center}
\begin{tabular}{ll}
\toprule
\textbf{qpos[1]} & \textbf{Physical meaning} \\
\midrule
$0$ rad & Pendulum upright (balance target, unstable equilibrium) \\
$\pm\pi$ rad & Pendulum hanging down (stable equilibrium) \\
\bottomrule
\end{tabular}
\end{center}

\noindent \textbf{qpos[0]} (shoulder) has no special reference -- $0$ corresponds to the arm's initial assembly position in SolidWorks.

%─────────────────────────────────────────────────────────────────────────────
\section{Gymnasium Environment}

The MuJoCo model is wrapped as a \texttt{gymnasium.Env} in \texttt{furuta\_env.py},
providing a standard reinforcement learning interface.

\subsection{Observation Space}

The 5-dimensional observation vector is:

\begin{equation}
\mathbf{o} = \bigl[\cos\theta,\ \sin\theta,\ \dot\theta,\ \phi,\ \dot\phi\bigr]
\end{equation}

\noindent where $\theta = \texttt{qpos[1]}$ (elbow/pendulum) and
$\phi = \texttt{qpos[0]}$ (shoulder/arm).
The angle $\theta$ is encoded as $(\cos\theta, \sin\theta)$ rather than the raw
angle to avoid the $\pm\pi$ discontinuity.

\subsection{Observation Bounds}

Velocity bounds are derived from the physical hardware, not set to infinity:

\begin{center}
\begin{tabular}{llrr}
\toprule
\textbf{Component} & \textbf{Derivation} & \textbf{Raw limit} & \textbf{Bound (1.5$\times$)} \\
\midrule
$\cos\theta,\ \sin\theta$ & trigonometric & $\pm 1$ & $\pm 1$ \\
$\dot\theta$ & free fall from upright (energy conservation) & 36.6 rad/s & $\pm 55$ rad/s \\
$\phi$ & shoulder hard limit $\pm135^\circ$ & $\pm2.356$ rad & $\pm2.5$ rad \\
$\dot\phi$ & terminal speed $= \tau_{\max}/b = 0.132/0.01$ & 13.2 rad/s & $\pm20$ rad/s \\
\bottomrule
\end{tabular}
\end{center}

\noindent A 1.5$\times$ safety margin covers transient states (e.g.\ energy coupling
during swing-up) that briefly exceed steady-state limits.

\subsection{Action Space}

A scalar $u \in [-1,\ 1]$ is passed directly to the shoulder motor actuator.
The MuJoCo motor \texttt{gear="0.132"} maps this linearly to
$[-0.132,\ 0.132]$ Nm.

\subsection{Reward}

The dense reward at each step is:

\begin{equation}
r = \cos(\theta) - 0.1u^2 - 0.5\max(0,\ |\phi|-2.094)^2
\end{equation}

\noindent This gives $r = +1$ when the pendulum is perfectly upright ($\theta = 0$)
with zero control effort and shoulder angle inside the soft limit.  The
control penalty discourages chattering, and the shoulder soft-limit penalty
activates beyond $\pm120^\circ$ before the hard $\pm135^\circ$ joint limit.

\subsection{Episode Structure}

\begin{center}
\begin{tabular}{ll}
\toprule
\textbf{Parameter} & \textbf{Value} \\
\midrule
MuJoCo timestep & 2 ms (default) \\
Frame skip & 5 \\
Control period & 10 ms \\
Episode length & 3000 steps (30 s) for current training runs \\
Termination & None (truncated at fixed episode length only) \\
Maximum return & $+3000$ for 30 s with perfect upright/zero effort \\
\bottomrule
\end{tabular}
\end{center}

\subsection{Reset}

Each episode resets the pendulum to the hanging-down position with small
uniform noise:

\begin{align}
\phi_0 &\sim \mathcal{U}(-0.05,\ 0.05) \text{ rad} \\
\theta_0 &= -\pi + \mathcal{U}(-0.05,\ 0.05) \text{ rad} \\
\dot\phi_0,\ \dot\theta_0 &\sim \mathcal{U}(-0.05,\ 0.05) \text{ rad/s}
\end{align}

\subsection{Environment Validation}

Eleven automated tests were run to verify correctness:

\begin{center}
\begin{tabular}{lll}
\toprule
\textbf{Test} & \textbf{What is checked} & \textbf{Result} \\
\midrule
Observation shape     & Shape is $(5,)$                               & PASS \\
Trig identity         & $\cos^2\theta + \sin^2\theta = 1$             & PASS \\
Reset position        & $\cos\theta \approx -1$ after reset           & PASS \\
Initial velocity      & $|\dot\theta|, |\dot\phi| < 0.2$ rad/s        & PASS \\
Upright reward        & $r = +1$ when $\theta = 0$                    & PASS \\
Hanging reward        & $r = -1$ when $\theta = \pm\pi$               & PASS \\
Episode length        & Fixed by \texttt{episode\_seconds}             & PASS \\
Termination flags     & \texttt{terminated=False, truncated=True}     & PASS \\
$\dot\theta$ bound    & Peak $|\dot\theta| < 55$ rad/s over 20 episodes & PASS \\
$\dot\phi$ bound      & Peak $|\dot\phi| < 20$ rad/s over 20 episodes  & PASS \\
Determinism           & Same seed produces identical trajectory        & PASS \\
\bottomrule
\end{tabular}
\end{center}

\noindent The bounds stress test applied alternating maximum torque ($u = \pm 1$) for
20 full episodes. Peak velocities observed: $|\dot\theta|_{\max} = 13.3$ rad/s
(76\% headroom) and $|\dot\phi|_{\max} = 14.0$ rad/s (30\% headroom).

%─────────────────────────────────────────────────────────────────────────────
\section{Reinforcement Learning: PPO Training}

\subsection{Algorithm}

Proximal Policy Optimisation (PPO) \cite{schulman2017proximal} is a model-free,
on-policy actor-critic algorithm. At each iteration it:
\begin{enumerate}
    \item Collects $N_{\text{envs}} \times n_{\text{steps}}$ transitions using the
          current policy.
    \item Estimates advantages using Generalised Advantage Estimation (GAE).
    \item Performs $n_{\text{epochs}}$ gradient updates with a clipped surrogate
          objective:
          \begin{equation}
          \mathcal{L}^{\text{CLIP}} =
          \mathbb{E}\Bigl[\min\bigl(r_t(\theta)\hat{A}_t,\;
          \text{clip}(r_t(\theta), 1-\epsilon, 1+\epsilon)\hat{A}_t\bigr)\Bigr]
          \end{equation}
          where $r_t(\theta) = \pi_\theta(a_t|s_t)/\pi_{\theta_{\text{old}}}(a_t|s_t)$
          and $\epsilon = 0.2$.
\end{enumerate}

\noindent The clipping prevents any single update from changing the policy too
drastically, ensuring stable learning without a trust-region constraint.

\subsection{Observation Normalisation}

Observations are normalised online using \texttt{VecNormalize}
(Stable-Baselines3). Running mean and variance are maintained across all
parallel environments. At each step:

\begin{equation}
\tilde{o}_i = \text{clip}\!\left(\frac{o_i - \mu_i}{\sigma_i},\ -10,\ +10\right)
\end{equation}

\noindent This is critical because the observation vector mixes dimensionless
trigonometric values ($\cos\theta, \sin\theta \in [-1,1]$) with velocities in
rad/s and an unbounded angle $\phi$. Without normalisation the value network
faces a $\sim\!100\times$ scale difference between dimensions.

\textbf{Important:} the normalisation statistics are saved at the exact moment
the best model checkpoint is written (via \texttt{callback\_on\_new\_best}),
ensuring the statistics are always synchronised with the corresponding model
weights.

\subsection{Hyperparameters}

\begin{center}
\begin{tabular}{lrl}
\toprule
\textbf{Parameter} & \textbf{Value} & \textbf{Rationale} \\
\midrule
\texttt{n\_envs}        & 4                   & parallel data collection \\
\texttt{n\_steps}       & 2048                & $\sim$20 s of data per env per update \\
\texttt{batch\_size}    & 64                  & standard mini-batch \\
\texttt{n\_epochs}      & 10                  & reuse each rollout 10 times \\
\texttt{gamma}          & 0.99                & long horizon, dense reward \\
\texttt{gae\_lambda}    & 0.95                & GAE bias-variance trade-off \\
\texttt{ent\_coef}      & 0.01                & entropy bonus for exploration during swing-up \\
\texttt{learning\_rate} & $3\times10^{-4}$    & Adam optimiser default \\
\texttt{clip\_range}    & 0.2                 & PPO clipping parameter $\epsilon$ \\
\texttt{total\_timesteps} & 3\,000\,000       & $\sim$24 min on CPU \\
\bottomrule
\end{tabular}
\end{center}

\subsection{Training Results}

\begin{center}
\begin{tabular}{lrr}
\toprule
\textbf{Milestone} & \textbf{Steps} & \textbf{Eval reward} \\
\midrule
Random policy (baseline)  & 0            & $-850$ \\
Begins learning swing-up  & $\sim$20k    & $+8$   \\
Consistent balance        & $\sim$370k   & $+933$ \\
Converged                 & 3\,000\,000  & $\mathbf{+946}$ \\
\bottomrule
\end{tabular}
\end{center}

\noindent The best model achieved an evaluation reward of $\mathbf{946 / 1000}$
across 5 deterministic episodes, corresponding to the pendulum remaining within
the upright zone for approximately 94.6\% of the 10-second episode.
Swing-up from the hanging position occurs in approximately 0.62 seconds.
Training speed on a standard laptop CPU is approximately 2100 steps/s,
requiring no GPU. The policy network is a two-layer MLP with 64 units per layer.

%─────────────────────────────────────────────────────────────────────────────
\section{Arm Joint Limit Constraint}

\subsection{Motivation}

On real hardware, continuous 360° arm rotation is not possible due to cable
wrap and mechanical stops. The shoulder joint is therefore limited to
$\pm135°$ (270° total travel), matching the physical hardware constraint.

\subsection{Model Change}

The joint limit is enforced in MuJoCo by adding \texttt{limited} and
\texttt{range} attributes to the shoulder joint:

\begin{lstlisting}[style=xml]
<joint name="shoulder" pos="0 0 0.02575" axis="0 0 1" damping="0.01"
       limited="true" range="-2.356 2.356"/>
\end{lstlisting}

\noindent MuJoCo enforces this as a hard constraint force at $\pm2.356$ rad
($\pm135°$). The observation space bound for $\phi$ is updated from $\pm\infty$
to $\pm2.5$ rad accordingly.

\subsection{Reward Shaping}

A soft penalty activates when the arm exceeds $\pm120°$ (2.094 rad), growing
quadratically toward the hard stop:

\begin{equation}
r = \underbrace{\cos(\theta)}_{\text{balance}}
  - \underbrace{0.5 \cdot \max(0,\; |\phi| - 2.094)^2}_{\text{joint limit penalty}}
\end{equation}

\noindent The penalty is zero throughout the central $\pm120°$ zone and reaches
only $-0.034$ at the hard stop ($\pm135°$), so it steers the policy away from
the wall without significantly interfering with the balance objective.

\subsection{Training Results with Joint Limit}

\begin{center}
\begin{tabular}{lrr}
\toprule
\textbf{Milestone} & \textbf{Steps} & \textbf{Eval reward} \\
\midrule
Random policy (baseline)  & 0            & $-850$ \\
Swing-up learned          & $\sim$20k    & $-73$  \\
Consistent balance        & $\sim$500k   & $+574$ \\
Temporary regression (adapting to limit) & $\sim$550k & $+150$ \\
Recovered and converged   & 3\,000\,000  & $\mathbf{+823}$ \\
\bottomrule
\end{tabular}
\end{center}

\noindent The joint limit costs approximately 123 reward points compared to the
unconstrained run (946 $\rightarrow$ 823), reflecting reduced arm freedom for
swing-up and balance corrections. The temporary regression at $\sim$550k steps
is characteristic of the policy reshaping its swing-up strategy from wide
rotational sweeps to tighter oscillations within the $\pm135°$ envelope.

\begin{center}
\begin{tabular}{lrr}
\toprule
\textbf{Configuration} & \textbf{Best eval reward} & \textbf{Episode reward} \\
\midrule
No joint limit  & 946 / 1000 & 947 / 1000 \\
$\pm135°$ limit & 823 / 1000 & 821 / 1000 \\
\bottomrule
\end{tabular}
\end{center}

\subsection{Residual Chattering}

Despite achieving 82\% upright time, the policy exhibits visible arm chattering
(rapid small oscillations) during balance. This is expected for three reasons:

\begin{enumerate}
    \item \textbf{No control cost in reward} -- the reward $\cos(\theta)$
          depends only on pendulum angle, not on how much the arm moves.
          The policy has no incentive to use smooth corrections over aggressive ones.
    \item \textbf{Joint limit pressure} -- with $\pm135°$ the arm is more
          constrained, leading to sharper corrections to stay within bounds.
    \item \textbf{Unstable equilibrium} -- any deviation from $\theta=0$
          is amplified by gravity, requiring continuous active correction.
\end{enumerate}

\noindent The standard fix is \textbf{control regularisation}: adding an
$-\lambda u^2$ term to penalise large torques, encouraging smooth corrections:

\begin{equation}
r = \cos(\theta) - \lambda u^2 - 0.5 \cdot \max(0,\; |\phi| - 2.094)^2
\end{equation}

\noindent with $\lambda = 0.1$ as the planned next training iteration.

%─────────────────────────────────────────────────────────────────────────────
\section{Control Regularisation}

\subsection{Reward Change}

A control cost term is added to penalise large torques and suppress the arm
chattering observed during balance:

\begin{equation}
r = \cos(\theta) - 0.1\,u^2 - 0.5 \cdot \max(0,\; |\phi| - 2.094)^2
\end{equation}

At maximum torque ($u = 1$) the penalty is $-0.1$ per step, small enough that
swing-up (which requires large torques briefly) is not significantly impaired,
while sustained chattering (many steps of moderate torque) becomes expensive.

\subsection{Training Setup}

\begin{center}
\begin{tabular}{ll}
\toprule
\textbf{Parameter} & \textbf{Value} \\
\midrule
Total steps     & 3\,000\,000 \\
Parallel envs   & 4 (\texttt{DummyVecEnv}) \\
Device          & CPU (\texttt{device="cpu"}) \\
Training time   & $\sim$23 min at 2133 fps \\
\bottomrule
\end{tabular}
\end{center}

\noindent \textbf{Note on GPU:} PyTorch was upgraded from the CPU-only build to
\texttt{torch 2.11.0+cu128} (CUDA 12.8) to support the NVIDIA RTX 5070
(Blackwell architecture). However, Stable-Baselines3's PPO with
\texttt{MlpPolicy} is faster on CPU than GPU for small networks: the
overhead of CPU$\leftrightarrow$GPU data transfer exceeds the benefit of
GPU matrix operations for a $[64, 64]$ MLP. Inference uses \texttt{device="cpu"}
throughout.

\subsection{Training Results}

\begin{center}
\begin{tabular}{lrr}
\toprule
\textbf{Configuration} & \textbf{Best eval reward} & \textbf{Recorded episode} \\
\midrule
$\pm135°$ limit, no ctrl cost & 823 / 1000 & 821 / 1000 \\
$\pm135°$ limit + $-0.1u^2$   & 799 / 1000 & \textbf{811 / 1000} \\
\bottomrule
\end{tabular}
\end{center}

\noindent The reward dropped $\sim$24 points due to the control cost term
being included in the evaluation score. Arm chattering was visibly reduced.
The 799 reward model is the baseline for all subsequent training.

%─────────────────────────────────────────────────────────────────────────────
\section{Domain Randomisation}

\subsection{Motivation}

The trained policy is optimised for the exact simulation parameters (mass,
damping, torque constant). Real hardware deviates from these values due to
manufacturing tolerance, voltage variation, and unmodelled friction.
Domain randomisation (DR) perturbs physical parameters each episode reset
so the policy learns control strategies that generalise across the parameter
distribution, making zero-shot sim-to-real transfer more likely to succeed.

\subsection{Sensor Architecture}

The two joints use fundamentally different sensing hardware, requiring separate
noise models:

\begin{center}
\begin{tabular}{lll}
\toprule
\textbf{Joint} & \textbf{Sensor} & \textbf{Characteristics} \\
\midrule
Shoulder (arm) & Motor built-in magnetic encoder & Pre-gearbox; high effective CPR \\
               &                                  & ($\sim$260 $\times$ 20 = 5200 eff. CPR); \\
               &                                  & gear backlash $\sim$1° at output \\
Elbow (pendulum) & AS5600 (12-bit, 4096 CPR)     & Direct on pivot; absolute; \\
                 &                                & must differentiate for velocity \\
\bottomrule
\end{tabular}
\end{center}

\subsection{Current DR Profile}

The current default DR profile in \texttt{train.py} is
\texttt{real\_ready\_stage2}.  It extends \texttt{real\_ready} with small
encoder zero-offset randomisation.  The profile still excludes elbow velocity
filtering, because filter lag reduced success rate.

\begin{center}
\begin{tabular}{lll}
\toprule
\textbf{Parameter} & \textbf{\texttt{real\_ready\_stage2} range} & \textbf{Reason} \\
\midrule
Arm/rod mass and inertia & $\pm5\%$ & CAD/scale uncertainty until measured \\
Shoulder damping         & $\pm5\%$ & Small gearbox/friction variation \\
Motor gear/torque scale  & $\pm5\%$ & Voltage and motor variation \\
Action delay             & 0--1 control steps & ESP-NOW/serial timing margin \\
Shoulder position noise  & $\sigma = 0.017$ rad & Gear backlash ($\sim$1 degree) \\
Shoulder velocity noise  & $\sigma = 0.050$ rad/s & Encoder velocity estimate \\
Elbow position noise     & $\sigma = 0.001$ rad & Small AS5600 absolute-angle noise \\
Elbow damping            & $[0,\ 0.00002]$ Nm s/rad & Tiny realistic pivot friction \\
Shoulder zero offset     & $\pm1^\circ$ & Encoder/calibration offset tolerance \\
Elbow zero offset        & $\pm0.5^\circ$ & AS5600 upright calibration tolerance \\
Elbow velocity filter    & Excluded & Adds phase lag; test separately \\
\bottomrule
\end{tabular}
\end{center}

\noindent Superseded older values are kept below only as historical context
from the first DR design; they should not be used for the current training run:

\begin{center}
\begin{tabular}{lll}
\toprule
\textbf{Parameter} & \textbf{Range} & \textbf{Physical source} \\
\midrule
Arm mass             & $\pm10\%$             & Manufacturing tolerance \\
Rod mass             & $\pm10\%$             & Manufacturing tolerance \\
Shoulder damping     & $\pm25\%$             & Bearing friction variation \\
Elbow damping        & $[0,\ 0.005]$         & Pivot friction (unmeasured) \\
Motor torque scale   & $\pm10\%$             & Voltage droop, motor variation \\
Action delay         & 0--1 steps            & ESP-NOW 1--2 ms jitter \\
Shoulder pos. noise  & $\sigma = 0.017$ rad  & Gear backlash ($\sim$1° at output) \\
Shoulder vel. noise  & $\sigma = 0.050$ rad/s & Pre-gear encoder, high CPR \\
Elbow pos. noise     & $\sigma = 0.003$ rad  & AS5600 $\sim$2 LSB \\
Elbow vel. filter    & $\alpha \in [0.6,\ 0.8]$ & Low-pass filter tuning on ESP32 \\
\bottomrule
\end{tabular}
\end{center}

\subsection{Per-Joint Velocity Estimation}

A key sim-to-real gap is velocity estimation. MuJoCo provides \texttt{qvel}
directly as an integrated state, but on real hardware:

\begin{itemize}
    \item \textbf{Shoulder}: the motor encoder (pre-gearbox) has high effective
          CPR, so pulse-counting gives clean velocity with small additive noise.
    \item \textbf{Elbow}: the AS5600 is an absolute position encoder; velocity
          must be estimated by differentiating consecutive readings:
          \begin{equation}
          \hat{\dot\theta}_t = \alpha \cdot \frac{\theta_t - \theta_{t-1}}{\Delta t}
          + (1 - \alpha)\,\hat{\dot\theta}_{t-1}
          \end{equation}
          This low-pass filter introduces phase lag.  Tests show that a broad
          filter range hurts robustness, while a narrow range can be tolerated.
          The current \texttt{real\_ready} policy therefore uses raw simulated
          \texttt{qvel[1]} during training and keeps elbow filtering as a
          hardware/firmware item to tune carefully.
\end{itemize}

\subsection{Training Status, June 4 2026}

All current training outputs are organized under \texttt{pendulum/runs/}.
The clean no-DR policy was trained for 30 s episodes and is used as the
warm-start source for DR fine-tuning.

\begin{center}
\begin{tabular}{llll}
\toprule
\textbf{Run} & \textbf{Folder} & \textbf{Best mean} & \textbf{Notes} \\
\midrule
No DR baseline & \texttt{runs/no\_dr/20260604\_163108} & 2886.57 / 3000 & Strong clean policy \\
Full DR & archived invalid & Failed & Collapsed under combined DR \\
\texttt{mass\_motor\_delay} & \texttt{20260604\_175615\_mass\_motor\_delay} & 2637.46 / 3000 & Passed mean, high variance \\
\texttt{real\_ready} & \texttt{20260604\_182638\_real\_ready} & 2771.58 / 3000 & High mean, poor worst-case robustness \\
\texttt{real\_ready\_stage2} & \texttt{20260604\_204533\_real\_ready\_stage2} & 2554.50 / 3000 & Current robust DR policy \\
\bottomrule
\end{tabular}
\end{center}

\noindent The \texttt{real\_ready} policy records a clean GIF score of
2896.0 / 3000 over 30 s, so nominal behavior is strong.  However, full
\texttt{real\_ready} randomized evaluation still has low worst-case episodes
(best-eval worst episode 188.80).  The Stage 2 policy lowers peak mean reward
slightly but improves worst-case robustness substantially: best mean
2554.50 / 3000, standard deviation 48.70, and worst episode 2439.84.

\subsection{No-DR vs Stage 2 Robustness}

To test sim-to-real robustness, the clean no-DR policy and the Stage 2 DR
policy were both evaluated under the same \texttt{real\_ready\_stage2}
randomized environment for 20 episodes.

\begin{center}
\begin{tabular}{lrrrrr}
\toprule
\textbf{Policy} & \textbf{Mean} & \textbf{Std} & \textbf{Worst} & \textbf{Best} & \textbf{Success} \\
\midrule
No DR baseline & 1729.64 & 1336.33 & 84.20 & 2893.91 & 60\% \\
Stage 2 DR & 2554.67 & 43.38 & 2483.71 & 2638.26 & 100\% \\
\bottomrule
\end{tabular}
\end{center}

\noindent This comparison applies Stage 2 randomized conditions to both
policies.  It shows that the no-DR policy is excellent in clean simulation but
fragile under realistic imperfections, while the Stage 2 policy is much more
consistent.

\begin{center}
\begin{tabular}{lrrrr}
\toprule
\textbf{Policy} & \textbf{Swing-up} & \textbf{Action RMS} & \textbf{Control Energy} & \textbf{Shoulder Range} \\
\midrule
No DR baseline & 0.96 s & 0.347 & 5.07 & 136.1$^\circ$ \\
Stage 2 DR & 1.02 s & 0.372 & 4.18 & 271.7$^\circ$ \\
\bottomrule
\end{tabular}
\end{center}

\subsection{Isolated DR Results}

\begin{center}
\begin{tabular}{llrl}
\toprule
\textbf{Profile} & \textbf{Run} & \textbf{Best mean} & \textbf{Conclusion} \\
\midrule
Mass/inertia & \texttt{20260604\_173556\_mass} & 2892.20 & Safe \\
Shoulder damping & \texttt{20260604\_173723\_shoulder\_damping} & 2866.00 & Safe at $\pm5\%$ \\
Motor & \texttt{20260604\_174320\_motor} & 2892.88 & Safe \\
Delay & \texttt{20260604\_174421\_delay} & 2909.24 & Safe after adaptation \\
Elbow damping & \texttt{20260604\_173826\_elbow\_damping} & -2141.73 & Fails \\
Sensor combined & \texttt{20260604\_174850\_sensor} & 2141.39 & Below target \\
\bottomrule
\end{tabular}
\end{center}

\subsection{Elbow Troubleshooting}

The elbow is currently the main sim-to-real risk.  The current XML has no
explicit elbow damping, but the old DR range randomly sampled up to
\texttt{0.003} Nm s/rad.  For the updated \texttt{rod\_2} model this is too
large.

\begin{center}
\begin{tabular}{rrl}
\toprule
\textbf{Fixed elbow damping} & \textbf{Mean reward} & \textbf{Result} \\
\midrule
0.00000 & 2896.99 & Stable \\
0.00002 & 2892.77 & Stable \\
0.00005 & 735.80  & Fails swing-up/balance \\
0.00010 & -269.17 & Fails \\
0.00050 & -3145.53 & Collapses \\
0.00300 & -3147.49 & Collapses \\
\bottomrule
\end{tabular}
\end{center}

\noindent The current rod mass is 0.01888 kg and the pivot-to-COM distance is
0.03603 m, giving a gravity torque scale of about 0.00667 Nm.  With damping
\texttt{c=0.003}, the damping torque at 10 rad/s is 0.030 Nm, roughly 4.5
times the gravity torque scale.  This explains why the old elbow damping DR
was not a mild perturbation.

\begin{center}
\begin{tabular}{rrrl}
\toprule
\textbf{Elbow sensor test} & \textbf{Best mean} & \textbf{Worst episode} & \textbf{Conclusion} \\
\midrule
Position noise only & 2864.88 & 2862.85 & Safe \\
Velocity noise only & 2891.53 & 2890.10 & Safe \\
Small elbow pos. noise & 2883.52 & 2860.66 & Safe \\
Broad elbow velocity filter & 2477.63 & 1981.34 & Fails target \\
Narrow elbow velocity filter & 2714.16 & 2215.98 & Passes mean, still variable \\
\bottomrule
\end{tabular}
\end{center}

\noindent Decision: keep small elbow position noise, tiny elbow damping, and
small encoder zero-offset DR in the current training profile.  Elbow damping
uses \texttt{uniform(0,\ 0.00002)} Nm s/rad.  Shoulder zero offset uses
$\pm1^\circ$ and elbow zero offset uses $\pm0.5^\circ$.  Elbow velocity
filtering remains excluded from \texttt{real\_ready\_stage2}; if the firmware
needs a filtered velocity estimate, test it separately with the existing
\texttt{sensor\_filter\_narrow} profile.

\subsection{Earlier Training Challenges}

\begin{enumerate}
    \item \textbf{Cold-start failure}: Training from scratch with full DR
          caused policy collapse by step 50k. The task (swing-up + balance
          under simultaneous noise, delay, and torque variation) was too hard
          for the policy to discover from random initialisation.

    \item \textbf{Fix -- warm-start}: Loading the best ctrl-cost model
          (799 reward) as the starting policy and fine-tuning with DR
          (\texttt{ent\_coef=0.05}, \texttt{learning\_rate=1e-4}) avoided the
          cold-start problem. The policy already knew how to swing up;
          DR fine-tuning only needed to teach robustness.

    \item \textbf{Observation mismatch}: The velocity filter was initially
          applied unconditionally in \texttt{\_get\_obs()}, changing the
          eval observation distribution from what the warm-start model expected.
          Fix: gate the filter inside \texttt{if self.\_domain\_rand}.

    \item \textbf{Delay too aggressive}: \texttt{\_max\_delay = 2} steps
          (20 ms over a 10 ms control loop) caused unstable fine-tuning.
          Reduced to \texttt{\_max\_delay = 1}.

    \item \textbf{Spinning failure mode}: One DR run learned to spin the
          pendulum continuously rather than balance, achieving a small positive
          reward by briefly passing through upright each revolution. Root cause:
          the ctrl cost alone was insufficient to discourage continuous rotation
          under DR conditions where weaker torque made each spin cheaper.
          A spin cost $-0.001\,\dot\theta^2$ was tested but interfered with
          swing-up (which requires transiently high $\dot\theta$). Resolved by
          tightening DR ranges so the ctrl cost was sufficient.

    \item \textbf{Plateau at $\sim$$-65$ to $-75$}: With tighter DR ranges
          and correct observation gating, the DR fine-tuned model plateaued
          around $-65$ eval reward and never consistently exceeded the initial
          warm-start score of $+155$ (measured before any gradient updates).
          This indicates curriculum DR is necessary for stable improvement.
\end{enumerate}

\subsection{Historical DR Model Result}

An earlier DR model (saved at step 10k, before fine-tuning degraded it) scored
$+155$ eval reward (clean conditions) and $+178$ in a recorded episode.
Visually: smooth arm motion, no spinning, but balance sustained for only
$\sim$1 second before the pendulum falls. The swing-up is reliable; balance
under perturbation remained the unsolved challenge at that stage.  This is
superseded by the 30 s \texttt{real\_ready} run documented above.

\subsection{Stopping Criteria}

To avoid wasting compute on plateaued runs, two automatic stopping callbacks
were added to \texttt{train.py}:

\begin{center}
\begin{tabular}{lll}
\toprule
\textbf{Callback} & \textbf{Trigger} & \textbf{Purpose} \\
\midrule
\texttt{StopTrainingOnRewardThreshold} & Eval reward $> 2550$ over 30 s & Stop early on success \\
\texttt{StopTrainingOnNoModelImprovement} & No new best after mature eval window & Stop on plateau \\
\bottomrule
\end{tabular}
\end{center}

%─────────────────────────────────────────────────────────────────────────────
\subsection{Curriculum Domain Randomisation}

\subsubsection{Motivation}

Jumping directly from a clean warm-start policy to full DR caused the policy to
collapse within 100--200k steps and never recover (eval reward $-65$ to $-75$,
vs.\ $+799$ clean). The fix is to \emph{ramp} DR gradually so the policy
adapts incrementally rather than facing the full perturbation all at once.

\subsubsection{Implementation}

A shared \texttt{schedule} dict (a Python object held in memory and passed to
all training environments via closure) stores a single scalar \texttt{progress}
$\in [0, 1]$. \texttt{CurriculumCallback} updates it every step:

\begin{lstlisting}[style=python]
class CurriculumCallback(BaseCallback):
    def _on_step(self) -> bool:
        self._schedule["progress"] = min(
            1.0, self.num_timesteps / self._curriculum_steps
        )
        return True
\end{lstlisting}

\noindent In \texttt{FurutaPendulumEnv.reset()}, all DR ranges scale linearly
with \texttt{progress} $p$:

\begin{lstlisting}[style=python]
if p > 0:
    mass_scale = uniform(1 - 0.05*p, 1 + 0.05*p)
    self.model.body_mass[...] *= mass_scale
    self.model.body_inertia[...] *= mass_scale
    self.model.dof_damping[0] = nom_shoulder_damp * uniform(1 - 0.05*p, 1 + 0.05*p)
    self.model.actuator_gear[0, 0] = nom_motor_gear * uniform(1 - 0.05*p, 1 + 0.05*p)
    delay = int(integers(0, round(p) + 1))   # 0 steps until p>0.5
    # real_ready_stage2 also adds shoulder/elbow encoder zero offsets
\end{lstlisting}

\noindent Observation noise in \texttt{\_get\_obs()} is also scaled by $p$.
For \texttt{real\_ready}, the sensor DR is shoulder position noise, shoulder
velocity noise, and small elbow position noise only.

\subsubsection{Safe Run Storage}

A dedicated \texttt{train\_clean.py} script trains the clean (no-DR) baseline
and saves all outputs to a timestamped directory under
\texttt{runs/no\_dr/}.  DR fine-tuning writes to a separate timestamped
directory under \texttt{runs/dr/}.  This keeps clean baselines, DR runs, logs,
normalization files, XML snapshots, and GIFs organized by experiment.

\begin{center}
\begin{tabular}{ll}
\toprule
\textbf{File} & \textbf{Written by} \\
\midrule
\texttt{runs/no\_dr/<timestamp>/best\_model.zip} & \texttt{train\_clean.py} \\
\texttt{runs/no\_dr/<timestamp>/vec\_normalize\_best.pkl} & \texttt{train\_clean.py} \\
\texttt{runs/dr/<timestamp>\_<profile>/best\_model.zip} & \texttt{train.py} \\
\texttt{runs/dr/<timestamp>\_<profile>/run\_config.json} & \texttt{train.py} \\
\bottomrule
\end{tabular}
\end{center}

\noindent \texttt{train.py} automatically loads the latest no-DR run as the
warm-start source and writes the selected DR profile into
\texttt{run\_config.json}.

\subsubsection{Clean Training Result}

The current clean baseline (no DR, fresh PPO from scratch, 30 s episodes)
reached \textbf{2886.57 / 3000} and is saved in
\texttt{runs/no\_dr/20260604\_163108/}.  The recorded nominal GIF achieved
approximately \textbf{2896 / 3000}.

\subsubsection{Historical Curriculum DR Training Results}

Earlier curriculum DR runs were performed before the current CAD/profile
cleanup.  They are retained as historical context only:

\begin{center}
\begin{tabular}{lllrr}
\toprule
\textbf{Run} & \textbf{ent\_coef} & \textbf{Curriculum} & \textbf{Best eval} & \textbf{Stopped at} \\
\midrule
DR run 1 & 0.05  & 2M steps & \textbf{808} (step 160k) & 1.0M (patience) \\
DR run 2 & 0.05  & 3M steps & 638 (step 10k)  & 2.5M (patience) \\
DR run 3 & 0.005 & 3M steps & 737 (step 220k) & 2.5M (patience) \\
\bottomrule
\end{tabular}
\end{center}

\noindent These runs used the older storage layout and older DR ranges.  They
are superseded by the current \texttt{runs/} experiments and the
\texttt{real\_ready} profile.

\subsubsection{Root Cause Analysis}

The warm-start policy peaks early because the DR training gradient updates
cause it to \emph{forget} precise balance timing before it learns DR robustness.
Two contributing factors:

\begin{enumerate}
    \item \textbf{High entropy coefficient} (\texttt{ent\_coef=0.05} vs.\
          \texttt{0.01} in clean training): the optimiser is penalised for
          being confident, pushing the policy to explore and abandon the
          learned balance behaviour. Reducing to \texttt{0.005} (run 3)
          improved early stability but did not solve the underlying issue.
    \item \textbf{Patience too short}: with \texttt{min\_evals=50}
          (original), the no-improvement check fired at 1M steps before the
          curriculum finished ramping. Extended to \texttt{min\_evals=150}
          and \texttt{max\_no\_improvement\_evals=100} in later runs (total
          steps increased to 4M), but the policy still never recovered its
          clean peak.
\end{enumerate}

\subsubsection{Next Direction}

The current best deployment candidate is the \texttt{real\_ready\_stage2} run
\texttt{runs/dr/20260604\_204533\_real\_ready\_stage2/}.  It crossed the
2550 / 3000 threshold and has strong worst-case randomized performance.  Before
hardware deployment, the highest-value measurements are elbow pivot friction
(free-swing decay), AS5600 velocity-estimator lag, motor response, and
end-to-end control delay.

\subsubsection{Encoder Zero Calibration}

The current elbow joint uses \texttt{ref="0"} and was verified in the MuJoCo
viewer: at \texttt{qpos[1]=0} the pendulum rod is visually upright.  The
policy's zero angle therefore means upright in simulation.

On the real robot, the AS5600 encoder must be calibrated so that it reads zero
when the rod is held upright. Without this, the policy sees a constant angular
offset and tries to balance at the wrong angle. The calibration procedure is:
hold the rod upright at power-on, record the raw encoder value, and subtract it
from all subsequent readings in firmware.

%─────────────────────────────────────────────────────────────────────────────
\section{Hardware Platform}

\subsection{Microcontroller}

The real-time controller is an \textbf{ESP32 (ESP-WROOM-32 / ESP32-S)},
running at 240 MHz (dual-core Xtensa LX6). It reads the AS5600 encoder and
drives the motor, targeting a 10 ms control loop to match the simulation
control period.

\subsection{Wireless Communication}

Real-time policy deployment uses \textbf{ESP-NOW}, Espressif's peer-to-peer
MAC-layer protocol. Compared to alternatives:

\begin{center}
\begin{tabular}{lrll}
\toprule
\textbf{Protocol} & \textbf{Latency} & \textbf{Jitter} & \textbf{Verdict} \\
\midrule
ESP-NOW   & 1--2 ms   & Very low & \textbf{Selected} \\
WiFi UDP  & 2--10 ms  & Medium   & Acceptable \\
WiFi TCP  & 5--20 ms  & High     & Too slow \\
BLE       & 10--50 ms & High     & Too slow \\
\bottomrule
\end{tabular}
\end{center}

\noindent The 1--2 ms ESP-NOW round-trip is well within the 10 ms control
period, making explicit latency randomisation during training unnecessary for
a first deployment attempt.

\subsection{System Architecture}

The trained policy runs on the PC where the full PyTorch model is available.
Commands are relayed to the pendulum over ESP-NOW via a bridge ESP32 connected
by USB:

\medskip
\begin{center}
\ttfamily
PC (PPO policy, Python)\\
$\updownarrow$ USB serial ($<$1 ms)\\
ESP32 bridge (serial $\leftrightarrow$ ESP-NOW relay)\\
$\updownarrow$ ESP-NOW ($\sim$1--2 ms)\\
ESP32 on pendulum (encoder read + motor drive)
\end{center}
\medskip

%─────────────────────────────────────────────────────────────────────────────
\section{Next Steps}

\subsection{Completed}
\begin{itemize}
    \item[$\checkmark$] Control regularisation ($-0.1u^2$), eval reward 799/1000
    \item[$\checkmark$] Domain randomisation infrastructure (per-joint noise, action delay, torque scale)
    \item[$\checkmark$] Warm-start fine-tuning strategy
    \item[$\checkmark$] Automatic stopping criteria (plateau + threshold)
    \item[$\checkmark$] CUDA 12.8 PyTorch install (RTX 5070 Blackwell)
    \item[$\checkmark$] Curriculum DR implementation (\texttt{CurriculumCallback}, scaled DR ranges)
    \item[$\checkmark$] Safe training pipeline with timestamped \texttt{runs/no\_dr/} and \texttt{runs/dr/} folders
    \item[$\checkmark$] Clean 30 s baseline retrained: 2886.57 / 3000 in \texttt{runs/no\_dr/20260604\_163108/}
    \item[$\checkmark$] Current DR candidate: \texttt{real\_ready\_stage2}, 2554.50 / 3000 mean and 2439.84 worst episode in \texttt{runs/dr/20260604\_204533\_real\_ready\_stage2/}
    \item[$\checkmark$] No-DR vs Stage 2 robustness report generated in \texttt{reports/stage2\_dr\_report\_20260604.md}
    \item[$\checkmark$] Elbow troubleshooting completed: damping and velocity filtering identified as weak points
    \item[$\checkmark$] Encoder zero convention verified in MuJoCo viewer (\texttt{ref="0"} correct)
\end{itemize}

\subsection{Immediate}
\begin{enumerate}
    \item \textbf{System identification} -- measure elbow pivot friction from
          free-swing decay, AS5600 velocity-estimator lag, motor response, and
          end-to-end latency.  Use these measurements to tighten DR ranges.
    \item \textbf{Disturbance testing} -- evaluate Stage 2 recovery from
          standardized elbow angle/velocity kicks after the pendulum is upright.
    \item \textbf{Robot deployment} -- test the current deployment candidate
          (\texttt{runs/dr/20260604\_204533\_real\_ready\_stage2/best\_model.zip} +
          \texttt{vec\_normalize\_best.pkl}) after AS5600 zero calibration.
\end{enumerate}

\subsection{Hardware Deployment}
\begin{enumerate}
    \setcounter{enumi}{3}
    \item \textbf{ESP32 firmware} -- implement AS5600 encoder read (I$^2$C),
          motor encoder read, motor PWM drive, and ESP-NOW packet handler.
          Implement low-pass velocity filter matching the simulated $\alpha$.
    \item \textbf{Bridge firmware} -- USB serial $\leftrightarrow$ ESP-NOW
          relay on the PC-side ESP32.
    \item \textbf{Real-hardware interface} -- replace \texttt{mj\_step} with
          real sensor reads and motor commands; deploy policy over ESP-NOW.
    \item \textbf{Sim-to-real debugging} -- systematic troubleshooting:
          verify observation signs, check velocity filter tuning, compare
          motor response to simulation.
\end{enumerate}

\begin{thebibliography}{9}
\bibitem{schulman2017proximal}
J.~Schulman, F.~Wolski, P.~Dhariwal, A.~Radford, and O.~Klimov,
``Proximal Policy Optimization Algorithms,''
\textit{arXiv preprint arXiv:1707.06347}, 2017.
\end{thebibliography}

\end{document}
